
%%
%% Part 8 - Appendices and Supplementary Material
%%
%% This final part compiles supplementary material referenced throughout
%% the report.  It includes placeholders for figures exported from
%% notebooks, comprehensive mathematical derivations, pseudocode for
%% algorithms, evaluation metric definitions, how-to-run instructions
%% and reproducibility guidelines.  These appendices preserve all
%% details from the original assignment while organising them in a
%% coherent structure.

\section{Notebook Figures}

During the execution of our experiments we generated numerous figures
illustrating intermediate results, training curves and sample grids.
Due to space constraints in the main text, these images are gathered
here.  Each figure is labelled with its original notebook cell
number for traceability.  Where applicable, placeholders indicate
where to insert the actual graphic files.

\begin{figure}[H]
  \centering
  \fbox{\parbox{0.9\linewidth}{\centering
    \textit{Figure Placeholder:}\newline
    Variance Scheduler – Linear Schedule (Cell 11, Output 2).}}
  \caption{Variance scheduler visualisation.}
\end{figure}

\begin{figure}[H]
  \centering
  \fbox{\parbox{0.9\linewidth}{\centering
    \textit{Figure Placeholder:}\newline
    Visualise Forward Diffusion Process (Cell 13, Output 2).}}
  \caption{Forward diffusion process on a sample digit.}
\end{figure}

\begin{figure}[H]
  \centering
  \fbox{\parbox{0.9\linewidth}{\centering
    \textit{Figure Placeholder:}\newline
    Additional dataset preparation visualisations.}}
  \caption{Dataset preparation steps.}
\end{figure}

The above placeholders correspond to figures that would be included
using \texttt{\includegraphics} or a similar command.  When
compiling this report with access to the image files, replace each
placeholder with the appropriate image file names.

\section{Mathematical Derivations}

\subsection{Marginal Distribution of the Forward Process}

We derive the closed-form expression for the marginal distribution
$q(x_t\mid x_0)$ used in diffusion models.  Starting with
\(x_t = \sqrt{\alpha_t}x_{t-1} + \sqrt{1-\alpha_t}\,\epsilon_{t-1}\),
we expand recursively and collect terms to arrive at
\(x_t = \sqrt{\bar{\alpha}_t}x_0 + \sqrt{1-\bar{\alpha}_t}\,\epsilon\).
This derivation is fully presented in Part~2; here we recall the
important steps and emphasise the independence of noise terms.

\subsection{Variational Lower Bound Decomposition}

The variational lower bound (VLB) on the log-likelihood decomposes as
\begin{equation}
  L_{\mathrm{VLB}} = L_0 + \sum_{t=1}^{T-1} L_t + L_T,
\end{equation}
where
\begin{align}
  L_0 &= \mathbb{E}_{q}\bigl[\log p_\theta(x_0\mid x_1)\bigr],\\
  L_t &= \mathbb{E}_{q}\bigl[D_{\text{KL}}\bigl(q(x_{t-1}\mid x_t,x_0)\| p_\theta(x_{t-1}\mid x_t)\bigr)\bigr],\\
  L_T &= \mathbb{E}_{q}\bigl[D_{\text{KL}}\bigl(q(x_T\mid x_0)\| p(x_T)\bigr)\bigr].
\end{align}
The KL terms encourage the learned reverse transitions to match the
true posteriors.  In practice they reduce to weighted MSE terms as
explained earlier.

\subsection{Score Matching and Flow Matching}

We show how the score matching objective in denoising score matching
reduces to the Flow Matching velocity regression.  For a Gaussian
perturbation kernel, the score satisfies
\(\nabla_{x_t}\log q(x_t\mid x_0) = (x_0 - x_t)/\sigma_t^2\).  By
choosing variance schedules such that $\sigma_t^2 = t(1-t)\sigma^2$ and
introducing a conditioning on a target sample $x_1$, we obtain a
constant velocity field equal to $x_1 - x_0$ along linear paths.  This
identifies the regression target used in Flow Matching.

\section{Algorithms and Pseudocode}

For completeness, we reproduce the pseudocode for training and
sampling algorithms discussed in previous chapters.  These snippets
can be used as templates when implementing the models from scratch.

\subsection{DDPM Training Loop}

\begin{lstlisting}[language=Python, caption=DDPM Training Pseudocode]
for each batch x0 in dataloader:
    t = random integer in [0, T)
    noise = standard normal
    xt = sqrt(alpha_bar_t[t]) * x0 + sqrt(1 - alpha_bar_t[t]) * noise
    noise_pred = unet(xt, t)
    loss = mse(noise_pred, noise)
    backpropagate and update parameters
\end{lstlisting}

\subsection{DDIM Sampling Loop}

\begin{lstlisting}[language=Python, caption=DDIM Sampling Pseudocode]
x = standard normal
for t in reversed(selected_timesteps):
    noise_pred = unet(x, t)
    alpha_bar_t = scheduler.alphas_cumprod[t]
    alpha_bar_prev = scheduler.alphas_cumprod[t_prev] if t_prev else 1
    pred_x0 = (x - sqrt(1 - alpha_bar_t) * noise_pred) / sqrt(alpha_bar_t)
    direction = sqrt(1 - alpha_bar_prev - eta**2 * (1 - alpha_bar_t/alpha_bar_prev)) * noise_pred
    noise = 0 if eta == 0 else standard normal
    x = sqrt(alpha_bar_prev) * pred_x0 + direction + eta * noise
\end{lstlisting}

\subsection{Flow Matching Training and Sampling}

\begin{lstlisting}[language=Python, caption=Flow Matching Training]
for each batch in dataloader:
    x1 = sample data
    x0 = sample noise
    t = random uniform in [0,1]
    xt = (1 - t) * x0 + t * x1
    target_velocity = x1 - x0
    pred_velocity = model(xt, t)
    loss = mse(pred_velocity, target_velocity)
    backpropagate and update parameters
\end{lstlisting}

\begin{lstlisting}[language=Python, caption=Flow Matching Sampling]
x = sample noise
for i in range(num_steps):
    t = i / num_steps
    v = model(x, t)
    x = x + v * (1/num_steps)
return x
\end{lstlisting}

\section{Evaluation Metrics Definitions}

For reference, we formally define the evaluation metrics used in
this report.

\paragraph{Sliced Wasserstein Distance.}  The SWD between two
probability measures $\mu$ and $\nu$ in $\mathbb{R}^d$ is
\begin{equation}
  \mathrm{SWD}(\mu,\nu) = \mathbb{E}_{\theta\sim \mathrm{Uniform}(S^{d-1})}
  \bigl[W_1\bigl(\langle \theta,\mu\rangle, \langle \theta,\nu\rangle\bigr)\bigr],
\end{equation}
where $W_1$ is the one-dimensional Wasserstein distance and
$\langle\theta,\cdot\rangle$ denotes projection onto the unit vector
$\theta$.

\paragraph{Autocorrelation Function (ACF).}  For a time series
$\{x_t\}$ with mean $\bar{x}$ and variance $\sigma^2$, the ACF at lag
$k$ is
\begin{equation}
  \rho(k) = \frac{\sum_{t=1}^{T-k} (x_t - \bar{x})(x_{t+k} - \bar{x})}{\sum_{t=1}^T (x_t - \bar{x})^2}.
\end{equation}

\paragraph{Skewness and Kurtosis.}  The skewness and kurtosis of a
random variable $X$ with mean $\mu$ and standard deviation $\sigma$ are
\begin{align}
  \gamma_1 &= \frac{\mathbb{E}[(X-\mu)^3]}{\sigma^3},\\
  \gamma_2 &= \frac{\mathbb{E}[(X-\mu)^4]}{\sigma^4} - 3.
\end{align}

\section{How to Run the Code}

To reproduce the experiments, follow these steps:

\begin{enumerate}[label=\arabic*., leftmargin=2em]
  \item Install the required Python packages listed in the
    accompanying \texttt{requirements.txt} using pip or conda.
  \item Navigate to the \texttt{code/} directory of the repository.
  \item Train the diffusion model by running \texttt{python run\_ddpm.py}.
  \item Train the Flow Matching model by running
    \texttt{python run\_flow\_matching.py}.
  \item (Optional) Run \texttt{python run\_stable\_diffusion.py} to
    fine-tune Stable Diffusion via DreamBooth.
  \item Compile this report using \LaTeX{} with
    \texttt{pdflatex main.tex}.
\end{enumerate}

Make sure you have a CUDA-capable GPU for efficient training.  Set the
random seed for reproducibility if needed.

\section{Assignment Specification Mapping}

Table~\ref{tab:assignmentmap} maps the content of this report to the
requirements of the original assignment.  Each row lists a requirement
and the corresponding section or part where it is satisfied.

\begin{table}[H]
  \centering
  \begin{tabular}{p{0.45\linewidth}p{0.45\linewidth}}
    \toprule
    Assignment Requirement & Location in Report \\
    \midrule
    Derive closed-form $q(x_t\mid x_0)$ & Part~2, Section 2.1 \\
    Discuss noise prediction vs mean & Part~2, Section 2.2 \\
    Explain VLB terms & Part~2, Section 2.3 \\
    Describe DDIM sampling & Part~2, Section 2.4 \\
    Explain classifier-free guidance & Part~2, Section 2.5 \\
    U-Net architecture details & Part~3, Section 3.1 \\
    Implement variance schedule & Part~3, Section 3.2 \\
    Implement training loop & Part~3, Section 3.3 \\
    Implement sampling & Part~3, Section 3.4 \\
    Present visual results & Part~3, Section 3.5 \\
    Describe Stable Diffusion and DreamBooth & Part~4 \\
    Vector field role & Part~5, Section 5.1 \\
    Linear path justification & Part~5, Section 5.2 \\
    Loss derivation & Part~5, Section 5.3 \\
    Implement Flow Matching & Part~6 \\
    Evaluation and metrics & Parts~6 and 7 \\
    Appendices, pseudocode and metrics & Part~8 \\
    \bottomrule
  \end{tabular}
  \caption{Mapping of assignment requirements to report sections.}
  \label{tab:assignmentmap}
\end{table}

\section{Originality Statement}

I hereby declare that this work is original and that all sources have
been properly cited.  The implementations are based on the referenced
papers, with appropriate modifications for the specific assignment
requirements.  All code was written by the author unless otherwise
noted.

\section{Contact Information}

For questions or feedback, please contact the assignment organisers at
\href{mailto:mollahoseini@ut.ac.ir}{mollahoseini@ut.ac.ir} or
\href{mailto:fatemehnadir@gmail.com}{fatemehnadir@gmail.com}.

\section{End of Report}

This concludes the appendices and the report as a whole.  Thank you
for reading.
